\chapter{Conclusiones y trabajo futuro}\label{cap:conclusiones}

\section{Resultados obtenidos}\label{sec:resultados-obtenidos}

El presente Trabajo Terminal diseñó la arquitectura de un sistema de gestión de expedientes penales para despachos jurídicos mexicanos, integrando seis servicios criptográficos fundamentales: integridad de documentos mediante SHA-256 \cite{nist2015fips1804}, confidencialidad en reposo mediante AES-256-GCM \cite{dworkin2007} con envelope encryption \cite{dworkin2012}, no repudio mediante firma digital RSA-3072 con SHA-256 y certificados X.509 emitidos por una CA interna, sellado de tiempo conforme a RFC 3161 \cite{adams2001} a través de Cincel como PSC autorizado bajo la NOM-151-SCFI-2016 \cite{nom151_2016}, autenticación robusta mediante Argon2id \cite{biryukov2021} para hashing de contraseñas y TOTP \cite{mraihi2011} como segundo factor, y revocación de certificados mediante CRL conforme a RFC 5280 \cite{cooper2008}.

El diseño se fundamentó en un análisis riguroso del marco legal mexicano aplicable a las transacciones entre particulares: el Código de Comercio \cite{ccom2018} (artículos 89, 89 bis, 93, 93 bis, 97 y 100), el Código Federal de Procedimientos Civiles \cite{cfpc2023} (artículo 210-A), el Código Nacional de Procedimientos Penales \cite{cnpp2021} (artículos 356 y 380) y la NOM-151-SCFI-2016. Este análisis permitió establecer que la firma electrónica simple, respaldada por un sello de tiempo de un PSC autorizado, constituye una estrategia probatoria válida y defendible en sede judicial, conforme a la cláusula de equivalencia funcional del artículo 100 del Código de Comercio.

La arquitectura hexagonal adoptada separa explícitamente el dominio de negocio de las dependencias externas mediante puertos e interfaces abstractas. Esta separación se manifiesta en dos decisiones de diseño clave: el módulo de firma acepta cualquier certificado X.509 RSA + SHA-256, independientemente de su emisor; y el módulo de sellado de tiempo opera contra cualquier proveedor que implemente RFC 3161. Ambas decisiones permiten sustituir la infraestructura criptográfica del prototipo (CA interna, Cincel) por infraestructura regulada (e.firma del SAT, PSC comercial) sin modificar la lógica de verificación.

\section{Contribuciones}\label{sec:contribuciones}

La contribución principal de este trabajo es el diseño de un puente documentado entre la teoría criptográfica contemporánea (NIST SP 800-57 \cite{barker2020}, RFC 9106 \cite{biryukov2021}, RFC 6238 \cite{mraihi2011}) y los requisitos específicos del marco legal mexicano para documentos electrónicos con valor probatorio. Cada decisión de diseño criptográfico se justificó con referencia tanto a los estándares internacionales que la sustentan como a los artículos del Código de Comercio y el CNPP que determinan su admisibilidad en sede judicial.

Como contribución adicional, el trabajo establece que la distinción entre firma electrónica simple y avanzada en el derecho mexicano no es de admisibilidad —ambas son admisibles conforme al artículo 89 del Código de Comercio— sino de carga de la prueba. Esta precisión, documentada con referencia a los artículos 97 y 100, permite argumentar que un prototipo académico con CA interna genera evidencia criptográfica válida sin necesidad de contratar un PSC para la emisión de certificados de firmante, siempre que el sello de tiempo provenga de un PSC autorizado.

\section{Limitaciones}\label{sec:limitaciones}

El presente trabajo presenta las siguientes limitaciones que deben considerarse al interpretar sus resultados. En primer lugar, el prototipo opera bajo una arquitectura de inquilino único (single-tenant): cada instancia sirve a un solo despacho. Un despliegue comercial requeriría una arquitectura multi-tenant con aislamiento de datos entre despachos, lo que introduce complejidad adicional en la gestión de claves y el control de acceso. En segundo lugar, la CA interna del prototipo emite certificados que constituyen firma electrónica simple, no avanzada; aunque esto es suficiente para la demostración académica, un despliegue productivo debería evaluar la migración a certificados emitidos por un PSC autorizado para obtener la presunción legal del artículo 97.

En tercer lugar, la confirmación empírica del algoritmo de firma de la TSA de Cincel queda pendiente: se presume sha256WithRSAEncryption por coherencia con la PKI mexicana, pero debe verificarse parseando un sello del sandbox con OpenSSL. En cuarto lugar, los parámetros de Argon2id (m = 46 MiB, t = 1, p = 1) corresponden a la configuración recomendada de OWASP \cite{owasp_pwstorage2024} pero deben calibrarse empíricamente en el hardware de despliegue para confirmar que el tiempo de hash se mantiene entre 0.5 y 1.0 segundo. Finalmente, la CRL del prototipo tiene una vigencia de siete días, lo que introduce una ventana teórica de exposición entre la revocación de un certificado y su detección por los verificadores; en un entorno productivo, esta ventana podría reducirse complementando la CRL con un responder OCSP.

\section{Trabajo futuro}\label{sec:trabajo-futuro}

Las líneas de trabajo futuro derivadas del presente diseño se organizan en tres horizontes. En el corto plazo, correspondiente a la fase de implementación del prototipo: (1) codificar los módulos criptográficos en Rust utilizando las bibliotecas ring, x509-cert y argon2 identificadas en el stack tecnológico; (2) obtener un sello de prueba del sandbox de Cincel y parsear el TimeStampToken con OpenSSL para confirmar el algoritmo de firma de la TSA; (3) implementar la suite de pruebas unitarias y de integración para los flujos de firma, cifrado, sellado de tiempo y verificación; y (4) calibrar empíricamente los parámetros de Argon2id en el hardware de despliegue.

En el mediano plazo, orientado a fortalecer la postura de seguridad del sistema: (1) integrar un módulo de seguridad de hardware (HSM) o un servicio de gestión de claves en la nube (KMS) para el almacenamiento de la KEK, eliminando la dependencia de variables de entorno; (2) implementar un responder OCSP que complemente la CRL para reducir la ventana de exposición ante revocaciones; (3) evaluar la migración del formato de sello de tiempo de .tsr independiente a CAdES-T o PAdES-T para ofrecer paquetes de verificación autocontenidos; y (4) incorporar WebAuthn/FIDO2 como opción progresiva de MFA para clientes con dispositivos compatibles, alcanzando el nivel AAL3 de NIST SP 800-63B \cite{nist2024}.

En el largo plazo, orientado a la viabilidad comercial: (1) rediseñar la arquitectura para soportar múltiples inquilinos (multi-tenant) con aislamiento criptográfico por despacho, donde cada tenant disponga de su propia KEK y su propio espacio de certificados; (2) migrar los certificados de firmante de la CA interna a un PSC autorizado para convertir la firma en avanzada conforme al artículo 97 del Código de Comercio \cite{ccom2018}, beneficiándose de la portabilidad arquitectónica diseñada desde el inicio; y (3) incorporar mecanismos de auditoría a nivel de sistema operativo y bitacorización inmutable (append-only) con firmas periódicas de la bitácora para fortalecer la cadena de custodia digital ante instancias judiciales.
